\section{Where Multi-Event Retrieval Fails}
\label{sec:diagnosis}

We first ask how far the pipeline is from a perfect answer at each layer, using the 40 annotated queries.

\paragraph{Video ranking.}
With the annotated event texts, the correct video is ranked first for 16 of 40 queries and lies beyond rank 20 for 7.
Five of these seven belong to the cooking series, where the median rank of the correct video is 9 against 1.5 for the other programs.

\paragraph{Alignment inside the correct video.}
We then give the system the correct video and decode the path from its score matrix.
Only 25\% of queries have every event placed correctly, and 43\% of events are correct (Table~\ref{tab:layers}).
Choosing for each event its highest-scoring frame, without any temporal constraint, gives 35\% event accuracy; DP raises it to 43\%, fixing 18 events and breaking 9 (paired 95\% bootstrap interval over queries $-2$ to $+18$ points).
This interval includes zero, so the observed gain is not evidence of a reliably positive effect across queries.
Among the 448 frames of a median video, a correct frame is ranked first for 35\% of events, within the top 20 for 76\%, and within the top 50 for 86\%.
Wrong frames are not near misses: the median distance between a wrongly selected frame and its interval is 45 seconds, and only 8\% lie within 5 seconds.
These distances indicate confusion between distinct moments rather than predominantly adjacent-frame errors.

\begin{table}[t]
\centering
\caption{Per-layer audit on 40 annotated queries (113 events) with the correct video given. Tolerance one second.}
\label{tab:layers}
\begin{tabular}{lc}
\toprule
Measurement & Value \\
\midrule
A correct frame ranked first within the video & 35\% of events \\
A correct frame within top 5 / 20 / 50 & 57\% / 76\% / 86\% \\
Event accuracy, best frame per event (no order) & 35\% \\
Event accuracy, DP path & 43\% \\
All events correct, DP path & 25\% \\
All events correct, best $\lambda$ per query (oracle) & 30\% \\
\bottomrule
\end{tabular}
\end{table}

\paragraph{The gap penalty.}
Whole-sequence correctness is sensitive to $\lambda$: 15\%, 17.5\%, 25\%, and 5\% for $\lambda=0$, $10^{-6}$, $10^{-5}$, and $10^{-4}$.
Even an oracle that picks the best $\lambda$ for every query from a grid of nine values only reaches 30\%, so a per-query choice of $\lambda$ has little headroom.
This bound is specific to the penalty family: because $\sum_i (t_{f_{i+1}}-t_{f_i}) = t_{f_M}-t_{f_1}$, the penalty constrains only the total span of the path, not how the gaps between events are distributed, and richer temporal models are not covered by it.
The value $10^{-5}$ was set during competition preparation; we report the sweep on the same queries as a sensitivity analysis, not as tuning.

\paragraph{Scoring ablations.}
The low within-video frame ranks motivated three training-free score interventions.
(i)~Normalizing each event row of $S_v$ (z-score, min--max), subtracting the best competing event at each frame, or smoothing over neighboring frames never exceeded the baseline 25\%.
(ii)~For video ranking we tested whether studio-wide appearance drowns episode-specific content by centering frame embeddings on the corpus, series, or video mean, or by cross-domain similarity local scaling (CSLS)~\cite{conneau2018wordtranslation} computed from 500 captions.
With SigLIP2 alone, raw scores give R@1 0.30 and R@10 0.50; every normalization was worse, and the cooking series degraded most (median rank 39 raw, 46--101 normalized).
(iii)~We evaluated a variant that keeps the top-$K$ frames per event and adds a directional transition score $\beta\,(v_b-v_a)^\top(q_{i+1}-q_i)/4$ between consecutive candidates, where $v$ are frame embeddings and $q$ event embeddings.
This term separates into $v_b^\top\Delta q$ and $-v_a^\top\Delta q$, so for a fixed event pair it is equivalent to reweighting the unary scores of the two frames; it is not a pairwise model of motion between them.
Its best setting ($K=16$, $\beta=3$) reached 25\% whole-sequence correctness and 46.9\% event accuracy, against 25\% and 43.4\% for the baseline.
Shuffling the transition scores lowered correctness to 5--17.5\%, which shows that the structure of the added scores matters but not that it captures motion; pruning to top-$K$ alone costs about five points that the term only recovers.
We also checked whether answers are duplicated across videos, as news segments are often rebroadcast: the highest SigLIP2 similarity between an answer frame and any other video was 0.967, and on inspection the closest matches were different episodes of the same format, not rebroadcasts.

Within the variants studied, the remaining error is therefore primarily associated with how frames of one video are scored against an event; other temporal models may still help. This motivates the two interactive mechanisms below: deciding early when the video itself is wrong, and letting the user repair individual events when it is right.
